\documentclass[11pt]{article}

\usepackage[margin=1in]{geometry}
\usepackage{amsmath}
\usepackage{amssymb}
\usepackage{booktabs}
\usepackage{hyperref}

\title{Quantitative Pathway for Coupling Atmosphere, Hydrosphere, Interior Thermal State, and Geochemistry}
\author{LIFEORIG working note}
\date{\today}

\begin{document}
\maketitle

\section{Purpose}

This document defines a quantitative pathway for deciding whether a planet has
a hydrosphere and, if it does, how the hydrosphere should provide boundary
conditions for ocean and geochemical calculations.

The intended sequence is:
\begin{enumerate}
    \item Solve the atmosphere.
    \item Use the atmosphere plus planet structure to evaluate liquid/ice water.
    \item Decide whether there is no hydrosphere, a surface ocean/pond, or an
    ice shell plus subsurface ocean.
    \item Use the hydrosphere state to define the pressure, temperature,
    composition, and mineral boundary conditions for PHREEQC/geochemistry.
\end{enumerate}

This note is written to be implementable. The details of specific ocean/ice
layer models can replace the simple formulas below, but the data flow should
remain the same.

\section{Atmospheric Inputs}

After the atmospheric solver converges, the hydrosphere module should receive:
\begin{align}
    z_i               &\quad \mathrm{altitude\ grid}\ [\mathrm{m}], \\
    T_i               &\quad \mathrm{temperature\ profile}\ [\mathrm{K}], \\
    P_i               &\quad \mathrm{pressure\ profile}\ [\mathrm{Pa}], \\
    n_{\mathrm{H_2O},i} &\quad \mathrm{H_2O\ number\ density\ profile}\ [\mathrm{m^{-3}}], \\
    T_s               &= T(z=0), \\
    P_s               &= P(z=0).
\end{align}

The atmosphere is the source of surface boundary conditions and of the
condensable vapor reservoir. If outgassing, delivery, or non-equilibrium
chemistry adds water to the system, it enters the atmosphere first and is then
seen by the hydrosphere through $n_{\mathrm{H_2O}}(z)$.

\section{Planetary Inputs}

The planet model should provide:
\begin{align}
    M_p      &\quad \mathrm{planet\ mass}\ [\mathrm{kg}], \\
    R_p      &\quad \mathrm{planet\ radius}\ [\mathrm{m}], \\
    g_s      &= \frac{G M_p}{R_p^2}, \\
    \rho_p   &= \frac{M_p}{(4/3)\pi R_p^3}.
\end{align}

The bulk density $\rho_p$ is not enough to determine composition uniquely, but
it can provide a first-order estimate of rock/ice fraction. This is especially
important for Europa-like bodies, where the ocean is not produced by
atmospheric condensation.

\section{Surface Liquid Water from Atmospheric Condensation}

\subsection{Saturation vapor pressure}

For each atmospheric layer, compute the saturation vapor pressure of water:
\begin{equation}
    p_{\mathrm{sat}}(T_i).
\end{equation}

A first implementation can use the Magnus formula:
\begin{equation}
    p_{\mathrm{sat}}(T) =
    610.94
    \exp\left[
        \frac{17.625\,T_C}{T_C + 243.04}
    \right]
    \quad [\mathrm{Pa}],
\end{equation}
where
\begin{equation}
    T_C = T[\mathrm{K}] - 273.15.
\end{equation}

This is a practical approximation near ordinary liquid-water conditions. More
accurate vapor pressure formulas can replace it later.

\subsection{Saturation number density}

The saturation number density of water vapor is
\begin{equation}
    n_{\mathrm{sat},i} =
    \frac{p_{\mathrm{sat}}(T_i)}{k_B T_i}.
\end{equation}

The excess water vapor number density is
\begin{equation}
    n_{\mathrm{excess},i}
    =
    \max\left[
        0,\,
        n_{\mathrm{H_2O},i} - n_{\mathrm{sat},i}
    \right].
\end{equation}

\subsection{Condensed water mass}

The atmospheric water vapor mass is
\begin{equation}
    M_{\mathrm{H_2O,atm}}
    =
    \int_0^{z_{\mathrm{top}}}
    n_{\mathrm{H_2O}}(z)\,
    m_{\mathrm{H_2O}}\,
    4\pi(R_p+z)^2\,dz .
\end{equation}

The saturation-capacity water mass is
\begin{equation}
    M_{\mathrm{H_2O,sat}}
    =
    \int_0^{z_{\mathrm{top}}}
    n_{\mathrm{sat}}(z)\,
    m_{\mathrm{H_2O}}\,
    4\pi(R_p+z)^2\,dz .
\end{equation}

The condensed water mass implied by the atmospheric solution is
\begin{equation}
    M_{\mathrm{cond}}
    =
    \int_0^{z_{\mathrm{top}}}
    n_{\mathrm{excess}}(z)\,
    m_{\mathrm{H_2O}}\,
    4\pi(R_p+z)^2\,dz .
\end{equation}

This is the surface-liquid reservoir directly implied by the atmosphere.

\subsection{Surface liquid condition}

Surface liquid water is allowed only if:
\begin{align}
    P_s &> P_{\mathrm{triple}} = 611.657\ \mathrm{Pa}, \\
    T_s &\geq 273.15\ \mathrm{K}, \\
    T_s &\leq T_{\mathrm{boil}}(P_s), \\
    M_{\mathrm{cond}} &> 0.
\end{align}

If these conditions are satisfied, the equivalent global surface-ocean depth is
\begin{equation}
    d_{\mathrm{surf}}
    =
    \frac{M_{\mathrm{cond}}}
    {4\pi R_p^2 \rho_w},
\end{equation}
where $\rho_w \approx 1000\ \mathrm{kg\,m^{-3}}$.

The bottom pressure of that surface ocean is
\begin{equation}
    P_{\mathrm{bottom,surf}}
    =
    P_s + \rho_w g_s d_{\mathrm{surf}}.
\end{equation}

\section{Interior Water/Ice Estimate from Bulk Density}

Atmospheric condensation cannot produce a Europa-like subsurface ocean. For
icy bodies, estimate a bulk ice/water fraction from the planet density.

\subsection{Two-component rock/ice model}

Assume a first-order two-component body:
\begin{equation}
    \rho_p =
    x_{\mathrm{ice}} \rho_{\mathrm{ice}}
    +
    (1-x_{\mathrm{ice}})\rho_{\mathrm{rock}},
\end{equation}
where $x_{\mathrm{ice}}$ is a volume fraction.

Solving:
\begin{equation}
    x_{\mathrm{ice}}
    =
    \frac{\rho_{\mathrm{rock}}-\rho_p}
    {\rho_{\mathrm{rock}}-\rho_{\mathrm{ice}}}.
\end{equation}

In implementation:
\begin{equation}
    x_{\mathrm{ice}} \leftarrow
    \mathrm{clip}(x_{\mathrm{ice}},0,1).
\end{equation}

Use first-pass constants:
\begin{align}
    \rho_{\mathrm{ice}}  &\approx 1000\ \mathrm{kg\,m^{-3}}, \\
    \rho_{\mathrm{rock}} &\approx 3300\ \mathrm{kg\,m^{-3}}.
\end{align}

The corresponding mass fraction is
\begin{equation}
    f_{\mathrm{ice,mass}}
    =
    \frac{x_{\mathrm{ice}}\rho_{\mathrm{ice}}}{\rho_p}.
\end{equation}

The total ice/water mass estimate is
\begin{equation}
    M_{\mathrm{ice/water}}
    =
    f_{\mathrm{ice,mass}} M_p.
\end{equation}

The equivalent global water-layer thickness is
\begin{equation}
    d_{\mathrm{water,eq}}
    =
    \frac{M_{\mathrm{ice/water}}}
    {4\pi R_p^2 \rho_w}.
\end{equation}

\subsection{Interpretation}

If $\rho_p > \rho_{\mathrm{rock}}$, the two-component model gives no ice/water
component. The body is treated as rocky/metal-rich in this first-pass model.

If $\rho_p < \rho_{\mathrm{ice}}$, the two-component model is insufficient:
porosity, low-density volatiles, or gas-rich structure may be required.

For Europa-like bodies, this model gives a physically motivated water/ice
reservoir that does not come from atmospheric vapor.

\section{Thermal Model for Ice Shell and Subsurface Ocean}

\subsection{Interior heat flux}

The interior model should provide an upward heat flux:
\begin{equation}
    q_{\mathrm{int}}\quad [\mathrm{W\,m^{-2}}].
\end{equation}

This can come from:
\begin{itemize}
    \item radiogenic heating,
    \item tidal heating,
    \item secular cooling,
    \item parameterized convection in the rocky mantle,
    \item a user-specified first-pass value.
\end{itemize}

\subsection{Conductive ice shell thickness}

For an icy body with surface temperature below melting, a first conductive ice
shell estimate is:
\begin{equation}
    q_{\mathrm{int}}
    =
    k_{\mathrm{ice}}
    \frac{T_m - T_s}{d_{\mathrm{ice}}},
\end{equation}
so
\begin{equation}
    d_{\mathrm{ice}}
    =
    k_{\mathrm{ice}}
    \frac{T_m - T_s}{q_{\mathrm{int}}}.
\end{equation}

Use first-pass constants:
\begin{align}
    k_{\mathrm{ice}} &\approx 2.5\ \mathrm{W\,m^{-1}\,K^{-1}}, \\
    T_m &\approx 273.15\ \mathrm{K}.
\end{align}

The model should allow replacement of $T_m$ by a pressure- and salinity-dependent
melting curve.

\subsection{Subsurface ocean condition}

A subsurface ocean exists in the first-pass model if:
\begin{align}
    M_{\mathrm{ice/water}} &> 0, \\
    T_s &< T_m, \\
    q_{\mathrm{int}} &> 0, \\
    d_{\mathrm{water,eq}} &> d_{\mathrm{ice}}.
\end{align}

Then:
\begin{align}
    d_{\mathrm{ocean,sub}} &=
    d_{\mathrm{water,eq}} - d_{\mathrm{ice}}, \\
    M_{\mathrm{ocean,sub}} &=
    4\pi R_p^2 \rho_w d_{\mathrm{ocean,sub}}.
\end{align}

The top pressure of the subsurface ocean is
\begin{equation}
    P_{\mathrm{ocean,top}}
    =
    P_s + \rho_{\mathrm{ice}} g_s d_{\mathrm{ice}}.
\end{equation}

The bottom pressure is
\begin{equation}
    P_{\mathrm{ocean,bottom}}
    =
    P_{\mathrm{ocean,top}}
    +
    \rho_w g_s d_{\mathrm{ocean,sub}}.
\end{equation}

\section{Ice Shell Convection Check}

The conductive shell estimate should be followed by a convection check. A
simple Rayleigh number is:
\begin{equation}
    Ra_{\mathrm{ice}}
    =
    \frac{
        \rho_{\mathrm{ice}} g_s \alpha_{\mathrm{ice}}
        \Delta T d_{\mathrm{ice}}^3
    }{
        \kappa_{\mathrm{ice}} \eta_{\mathrm{ice}}
    },
\end{equation}
where
\begin{align}
    \Delta T &= T_m - T_s, \\
    \alpha_{\mathrm{ice}} &\quad \mathrm{thermal\ expansivity}, \\
    \kappa_{\mathrm{ice}} &\quad \mathrm{thermal\ diffusivity}, \\
    \eta_{\mathrm{ice}} &\quad \mathrm{effective\ ice\ viscosity}.
\end{align}

If
\begin{equation}
    Ra_{\mathrm{ice}} > Ra_{\mathrm{crit}},
\end{equation}
then a purely conductive shell is not self-consistent. A first critical value is
\begin{equation}
    Ra_{\mathrm{crit}} \sim 10^3,
\end{equation}
but stagnant-lid ice convection uses more detailed scalings.

In code, this can initially be stored as a diagnostic:
\begin{equation}
    \mathrm{ice\_shell\_regime}
    =
    \begin{cases}
    \mathrm{conductive}, & Ra_{\mathrm{ice}} \leq Ra_{\mathrm{crit}}, \\
    \mathrm{convective}, & Ra_{\mathrm{ice}} > Ra_{\mathrm{crit}}.
    \end{cases}
\end{equation}

\section{Rocky Interior Thermal Evolution}

For rocky planets or rocky cores, use parameterized convection rather than full
2D/3D convection as the first implementation.

\subsection{Energy balance}

Let $T_m$ be a representative mantle/core temperature. A simple thermal
evolution equation is:
\begin{equation}
    \rho_m C_p V_m \frac{dT_m}{dt}
    =
    H(t) - Q_{\mathrm{out}},
\end{equation}
where:
\begin{align}
    \rho_m &\quad \mathrm{mantle/core\ density}, \\
    C_p &\quad \mathrm{heat\ capacity}, \\
    V_m &\quad \mathrm{convecting\ volume}, \\
    H(t) &\quad \mathrm{internal\ heat\ production}, \\
    Q_{\mathrm{out}} &\quad \mathrm{heat\ loss}.
\end{align}

The heat loss is also the lower-boundary heat source for the atmosphere and
surface energy balance. If $Q_{\mathrm{out}}$ is a total power, then the
surface internal flux is:
\begin{equation}
    F_{\mathrm{int}}
    =
    \frac{Q_{\mathrm{out}}}{4\pi R_p^2}.
\end{equation}
If the interior model directly computes a heat flux $q_{\mathrm{int}}$, then:
\begin{equation}
    F_{\mathrm{int}} = q_{\mathrm{int}}.
\end{equation}

\subsection{Internal heat production}

The total internal heat production is decomposed as:
\begin{equation}
    H(t)
    =
    H_{\mathrm{rad}}(t)
    +
    H_{\mathrm{tidal}}(t)
    +
    H_{\mathrm{grav}}(t).
\end{equation}

\subsubsection{Radiogenic heating}

For a detailed isotope model:
\begin{equation}
    H_{\mathrm{rad}}(t)
    =
    \sum_i H_{i,0}\exp(-\lambda_i t),
\end{equation}
where
\begin{equation}
    \lambda_i = \frac{\ln 2}{t_{1/2,i}}.
\end{equation}
The main long-lived isotopes in rocky interiors are:
\begin{equation}
    ^{238}\mathrm{U},\quad
    ^{235}\mathrm{U},\quad
    ^{232}\mathrm{Th},\quad
    ^{40}\mathrm{K}.
\end{equation}

For a first implementation, a lumped effective decay law is acceptable:
\begin{equation}
    H_{\mathrm{rad}}(t)
    =
    H_{\mathrm{rad},0}\exp\left(-\frac{t}{\tau_{\mathrm{rad}}}\right).
\end{equation}

\subsubsection{Tidal heating}

For a synchronously rotating satellite or close-in planet with eccentricity
tides, a first-order total tidal power is:
\begin{equation}
    H_{\mathrm{tidal}}
    =
    \frac{21}{2}
    \frac{k_2}{Q}
    \frac{
        G M_\star^2 R_p^5 n e^2
    }{
        a^6
    },
\end{equation}
where:
\begin{align}
    k_2 &\quad \mathrm{second\ Love\ number}, \\
    Q &\quad \mathrm{tidal\ quality\ factor}, \\
    M_\star &\quad \mathrm{host/primary\ mass}, \\
    R_p &\quad \mathrm{planet/moon\ radius}, \\
    n &\quad \mathrm{mean\ motion}, \\
    e &\quad \mathrm{orbital\ eccentricity}, \\
    a &\quad \mathrm{orbital\ semi-major\ axis}.
\end{align}
The mean motion is:
\begin{equation}
    n =
    \sqrt{\frac{G M_\star}{a^3}}.
\end{equation}

The corresponding globally averaged tidal heat flux is:
\begin{equation}
    F_{\mathrm{tidal}}
    =
    \frac{H_{\mathrm{tidal}}}{4\pi R_p^2}.
\end{equation}

This term is essential for Europa-like cases, where the atmosphere does not
control the existence of the subsurface ocean but tidal heat can maintain a
liquid layer below an ice shell.

\subsubsection{Gravitational and secular heat}

Part of gravitational/secular cooling is already represented by the thermal
inertia term on the left-hand side of the energy equation:
\begin{equation}
    \rho_m C_p V_m \frac{dT_m}{dt}.
\end{equation}
If an additional early-time gravitational energy source is desired, use a
parameterized form:
\begin{equation}
    H_{\mathrm{grav}}(t)
    =
    H_{\mathrm{grav},0}
    \exp\left(-\frac{t}{\tau_{\mathrm{grav}}}\right).
\end{equation}
This can represent early differentiation, contraction, or accretional residual
heat in a first-pass model.

\subsection{Nusselt-Rayleigh scaling}

The convective heat flux can be parameterized as:
\begin{equation}
    Nu = a Ra^\beta.
\end{equation}

Then:
\begin{equation}
    q_{\mathrm{conv}}
    =
    Nu\,
    \frac{k_m(T_m-T_s)}{D_m},
\end{equation}
and
\begin{equation}
    Q_{\mathrm{out}} =
    4\pi R_p^2 q_{\mathrm{conv}}.
\end{equation}

The Rayleigh number is:
\begin{equation}
    Ra =
    \frac{
        \rho_m g_s \alpha_m (T_m-T_s) D_m^3
    }{
        \kappa_m \eta_m
    }.
\end{equation}

First-pass values often use $\beta \approx 1/3$, but stagnant-lid convection
requires care because viscosity is strongly temperature dependent.

\section{Hydrosphere State Classification}

After computing surface condensation and interior ice/water state, classify the
planet:

\begin{enumerate}
    \item \textbf{No hydrosphere}
    \begin{align}
        M_{\mathrm{cond}} &= 0, \\
        M_{\mathrm{ice/water}} &= 0
    \end{align}
    or liquid/ice conditions fail.

    \item \textbf{Surface ocean or pond}
    \begin{align}
        M_{\mathrm{cond}} &> 0, \\
        P_s &> P_{\mathrm{triple}}, \\
        273.15\ \mathrm{K} \leq T_s &\leq T_{\mathrm{boil}}(P_s).
    \end{align}

    \item \textbf{Ice shell with subsurface ocean}
    \begin{align}
        T_s &< T_m, \\
        d_{\mathrm{water,eq}} &> d_{\mathrm{ice}}.
    \end{align}

    \item \textbf{Ice-only hydrosphere}
    \begin{align}
        T_s &< T_m, \\
        d_{\mathrm{water,eq}} &\leq d_{\mathrm{ice}}.
    \end{align}
\end{enumerate}

\section{Geochemical Boundary Conditions}

PHREEQC/geochemistry should run only after the hydrosphere state is known.

\subsection{Surface pond or surface ocean}

For a surface pond/ocean:
\begin{align}
    T_{\mathrm{geo}} &= T_s, \\
    P_{\mathrm{geo}} &= P_s + \rho_w g_s d_{\mathrm{surf}}, \\
    \mathrm{composition} &= \mathrm{planet.hydro[ocean\_composition]}, \\
    \mathrm{minerals} &= \mathrm{surface\ rock\ preset}.
\end{align}

This corresponds to surface aqueous geochemistry, weathering, and volcanic-rock
wetting environments.

\subsection{Volcanic rock environment}

For volcanic-rock chemistry:
\begin{align}
    T_{\mathrm{geo}} &= T_s \quad \mathrm{or\ a\ volcanic\ surface\ temperature}, \\
    P_{\mathrm{geo}} &= P_s, \\
    \mathrm{fluid} &= \mathrm{condensed\ water/pond\ composition}, \\
    \mathrm{minerals} &= \mathrm{basaltic/volcanic\ rock\ assemblage}.
\end{align}

This environment is relevant when surface liquid water exists but does not cover
all rocky surface. The model should track available exposed rock fraction if the
surface ocean depth is shallow.

\subsection{Hydrothermal vent}

For hydrothermal vent chemistry:
\begin{align}
    T_{\mathrm{geo}} &= T_{\mathrm{reaction\ depth}}, \\
    P_{\mathrm{geo}} &= P_{\mathrm{ocean,bottom}}, \\
    \mathrm{fluid} &= \mathrm{ocean\ composition}, \\
    \mathrm{minerals} &= \mathrm{rocky\ core/basalt/ultramafic\ assemblage}.
\end{align}

The reaction temperature should come from the interior thermal model. A first
implementation can choose a reaction depth or a temperature window, for example:
\begin{equation}
    T_{\mathrm{reaction}} \in [350, 650]\ \mathrm{K}.
\end{equation}

The hydrothermal flux to the ocean can later be written:
\begin{equation}
    F_j =
    \dot{M}_{\mathrm{fluid}}
    \left(C_{j,\mathrm{vent}} - C_{j,\mathrm{ocean}}\right),
\end{equation}
where $\dot{M}_{\mathrm{fluid}}$ is a water circulation rate and $C_j$ is the
concentration of species $j$.

\section{Suggested Code Interface}

The main implementation target should be:
\begin{verbatim}
compute_hydrosphere_state(
    planet_mass,
    planet_radius,
    surface_temperature,
    surface_pressure,
    gravity,
    altitude,
    temperature_profile,
    h2o_number_density,
    internal_heat_flux,
    rho_rock=3300 kg/m3,
    rho_ice=1000 kg/m3,
)
\end{verbatim}

It should return a state object containing:
\begin{verbatim}
has_surface_ocean
surface_ocean_mass
surface_ocean_depth
has_subsurface_ocean
ice_shell_thickness
subsurface_ocean_mass
subsurface_ocean_depth
ice_shell_regime
ocean_top_pressure
ocean_bottom_pressure
hydrosphere_class
\end{verbatim}

\section{References}

\begin{enumerate}
    \item Schubert, G., Turcotte, D. L., and Olson, P. \textit{Mantle
    Convection in the Earth and Planets}. Cambridge University Press.
    \url{https://www.cambridge.org/core/books/mantle-convection-in-the-earth-and-planets/}

    \item Thiriet, M., Breuer, D., Michaut, C., and Plesa, A.-C. ``Scaling laws
    of convection for cooling planets in a stagnant lid regime.'' \textit{Physics
    of the Earth and Planetary Interiors}, 286, 138--153, 2019.
    \url{https://doi.org/10.1016/j.pepi.2018.11.003}

    \item Korenaga, J. ``Scaling of stagnant-lid convection with Arrhenius
    rheology and the effects of mantle melting.'' \textit{Geophysical Journal
    International}, 179, 154--170, 2009.
    \url{https://academic.oup.com/gji/article/179/1/154/2007792}

    \item Mitri, G. and Showman, A. P. ``Thermal convection in ice-I shells of
    Titan and Enceladus.'' \textit{Icarus}, 2008.
    \url{https://experts.arizona.edu/en/publications/thermal-convection-in-ice-i-shells-of-titan-and-enceladus}

    \item Tobie, G. and collaborators. ``Tidal Deformation and Dissipation
    Processes in Icy Worlds.'' Review article.
    \url{https://pmc.ncbi.nlm.nih.gov/articles/PMC11739232/}

    \item Kang, W. and Flierl, G. ``Europa's Coupled Ice-Ocean System: Temporal
    Evolution of a Pure Ice Shell.'' arXiv preprint.
    \url{https://arxiv.org/abs/2309.16821}

    \item Vance, S. D., Hand, K. P., and Pappalardo, R. T. ``Geophysical
    controls of chemical disequilibria in Europa.'' Related literature on
    subsurface ocean chemistry and exchange processes. For a broad review, see:
    \url{https://experts.azregents.edu/en/publications/subsurface-water-oceans-on-icy-satellites-chemical-composition-an/}

    \item Reactive transport modeling background:
    \url{https://en.wikipedia.org/wiki/Reactive_transport_modeling_in_porous_media}

    \item Serpentinization background and hydrothermal H2 production:
    \url{https://en.wikipedia.org/wiki/Serpentinization}
\end{enumerate}

\end{document}
